\section{Removing the clock: one decoder at every stopping length}\label{sec:auto}
We now change the query-time specification explicitly. The commands still
come from the same external alphabet, and there is still only one terminal
query. It may, however, arrive after any number $0\le n\le N$ of commands,
and the answer must have mean $e_j^T R_\alpha^{\sum c_t}x/10$ at that
length. The machine does not receive $n$ as a separate input.

\begin{definition}[Anytime autonomous realization]\label{def:anytime}
An anytime autonomous machine has one finite set of available states,
fixed command matrices $T_0,T_1$, fixed initialization rows, and one
fixed decoder for each query. The same rows work for every command length
from zero through $N$. Let $A_N(\alpha)$ be the least total number of
states, counting any states used to retain the emitted answer.
The device may depend on the prescribed maximum $N$; its rows may not
depend on the actual stopping length. Tables and exact atomic sampling
remain free, so this is not a succinct bit-space model.
\end{definition}

\begin{lemma}[Finite positive peripheral spectrum]\label{lem:peripheral}
Every eigenvalue of modulus one of a finite row-stochastic matrix is a
root of unity.
\end{lemma}
\begin{proof}
Let $Tv=\lambda v$ and choose $i$ with $|v_i|$ maximal and positive.
Equality in the triangle inequality for
$\lambda v_i=\sum_jT_{ij}v_j$ forces
$v_j=\lambda v_i$ for every positive successor of $i$.
All these successors have the same maximal modulus. Iteration follows
a positive directed path, which eventually contains a cycle of length
$l\ge1$. Along that cycle, $v_j=\lambda^l v_j$ with $v_j\ne0$.
Thus $\lambda^l=1$. This is the familiar peripheral-spectrum constraint
for positive realizations; see \cite{BF03,BF04}.
\end{proof}

\begin{theorem}[Linear autonomous state complexity]\label{thm:autonomous}
For any infinite-order planar rotation and every $N\ge1$,
\begin{equation}\label{eq:anytime}
                         N+1\le A_N(\alpha)\le3(N+1)+2.
\end{equation}
The upper machine has rational rows when $R_\alpha$ is rational.
Thus $A_N(\alpha)=\Theta(N)$, whereas the fixed-length vector-state
width is $\Theta_\alpha(N^{1/3})$ for bounded-type angles.
\end{theorem}
\begin{proof}
Suppose a machine has $K\le N$ total states. Fix one nonzero seed and
use only active commands. Let $d=d_1+i d_2$ be the complex column formed
from its two fixed terminal means and let $z\ne0$ be the initial
complex predictive vector. With $T=T_1$ and $\zeta=e^{2\pi i\alpha}$,
correctness says
\[
                         E T^j d=z\zeta^j\quad(0\le j\le N).
\]
Write the characteristic polynomial of $T$ as
$\chi(X)=\sum_{j=0}^K a_jX^j$. Cayley--Hamilton and the first $K+1$
required lengths give
\[
             0=E\chi(T)d=z\sum_{j=0}^K a_j\zeta^j=z\chi(\zeta).
\]
Thus $\zeta$ is an eigenvalue of $T$, contradicting
Lemma~\ref{lem:peripheral}. Hence $K\ge N+1$.
The argument needs only the finite length interval, not an infinite
stationary specification or a minimality assumption.

For the upper bound use the labels $(s,k)$, $1\le s\le3$, $0\le k\le N$,
from the triangle-counter construction in Theorem~\ref{thm:rational},
now all available in one state set. The idle matrix is the identity;
the active matrix increments $k$, with arbitrary normalized rows after
$k=N$. The decoder at $(s,k)$ reads $R_\alpha^k v_s$ and has no epoch
argument. Its norm is less than one. Initialization is the seed's
triangle decomposition. The same matrices and decoder therefore answer
at every permitted length. Add two absorbing answer states if the output
must persist after its emission. This proves the count and rationality.
\end{proof}

Equation~\eqref{eq:anytime} completes the proof of
Theorem~\ref{thm:main}. It does not lower-bound an autonomous device
required to answer only at one fixed length and allowed to be incorrect
at all earlier lengths. That is another problem. In particular, the
resonant construction in Section~\ref{sec:resonance} already has
stationary command rows, but its decoder compensates for exactly $N$
normalizing contractions. The strengthened anytime requirement is what
prevents that compensation from concealing the elapsed length.
The finite-horizon comparison is a concrete use of classical spectral
constraints, not a claim to discover the stationary rotation obstruction.
